\documentclass[tikz,border=2mm]{standalone}
\usetikzlibrary{math}

\makeatletter
\newif\iffounddigit

% 1. 递归检查每一个字符
\def\scan@for@digits#1{%
  \ifx#1\@empty
    % 扫描结束
  \else
    % 针对 0-9 每一个数字进行硬比对
    \if#10\global\founddigittrue\fi
    \if#11\global\founddigittrue\fi
    \if#12\global\founddigittrue\fi
    \if#13\global\founddigittrue\fi
    \if#14\global\founddigittrue\fi
    \if#15\global\founddigittrue\fi
    \if#16\global\founddigittrue\fi
    \if#17\global\founddigittrue\fi
    \if#18\global\founddigittrue\fi
    \if#19\global\founddigittrue\fi
    % 继续处理下一个 token
    \expandafter\scan@for@digits
  \fi
}

% 2. 预处理：将内容转为字符串并剔除“\”的影响
\newcommand{\checkfornumbers}[1]{%
  \global\founddigitfalse
  % 将参数转为纯文本字符串
  \edef\temp@str{\detokenize{#1}}%
  % 启动递归扫描
  \expandafter\scan@for@digits\temp@str\@empty
}

\tikzset{
  define/@circle-point/.code args={#1,#2}{%
  \message{^^J[PASS] \detokenize{#2} is a Point^^J}%
},
define/@circle-radius/.code args={#1,#2}{%
  \message{^^J[MATCH] \detokenize{#2} contains Numbers/Expression^^J}%
},
  define/circle/.code args={#1,#2}{%
    % \message{^^J#2}
    \checkfornumbers{#2}%
    \iffounddigit
      \pgfkeys{/tikz/define/@circle-radius={#1,#2}}%
    \else
      \pgfkeys{/tikz/define/@circle-point={#1,#2}}%
    \fi
  },
  doublee/.code args={#1,#2}{
    #2 = 2* #1;
  },
}
\makeatother

\begin{document}
\begin{tikzpicture}
  \tikzmath{
    \O{1} = 1; \O{2} = 1;
    \P{1} = 2; \P{2} = 3;
    \radius = 5;
  }

  % 情况 A：只有 \P，不含数字 0-9，应该识别为 Point
  \tikzset{define/circle={\O,\P}}

  % 情况 B：含数字 1，应该识别为 Radius
  \tikzset{define/circle={\O,\P{1}}}

  % 情况 C：含常数，应该识别为 Radius
  \tikzset{define/circle={\O, 5cm}}

  \tikzset{
    doublee={\radius,\d},
    define/circle={\O,\d}}

\end{tikzpicture}
\end{document}

% 这是一个非常经典的问题：在 LaTeX 中，宏的展开（Expansion）时机决定了逻辑判断的结果。

% 为什么最后一个识别为 Point？
% 在你的代码中，最后一个测试用例是：
% \tikzset{double={\radius,\r}, define/circle={\O, \r}}

% 赋值阶段：double 样式虽然计算了结果，但 \r 此时只是一个存储了数值的宏（或者在 tikzmath 作用域内的变量）。

% 检测阶段：当你调用 define/circle={\O, \r} 时，\checkfornumbers 接收到的参数 #2 是 \r 这个原始记号（Token）。

% 字符串化：\detokenize{#1} 将 \r 转换成了字符串 "\r "（反斜杠 + 字母 r + 空格）。

% 扫描结果：你的递归扫描函数 \scan@for@digits 只检查字符 0-9。因为字符串 "\r " 中不包含任何数字字符，所以 \iffounddigit 返回 false，最终被误判为 Point。